% GENERATED FILE. Edit canonical lessons, then run npm run generate:books.
% canonical-source-hash: fnv1a64:57893773e4a56bbb
% canonical-lessons: TA-C518-niral, TA-C518-kuriyitu, TA-C518-kappuppirati, TA-C518-ropo, TA-C518-ceyarkai

\chapter{Code and Robots}
\label{ch:ta-code-and-robots}
\hlchaptermodality{518}

\begin{chapteropening}
\textbf{By the end of this chapter:} \emph{I can talk about programs, codes and robots.}

Complete the last lesson of chapter 518: I can talk about programs, codes and robots.
\end{chapteropening}

\section[\texorpdfstring{\ta{நிரல்} (niral)}{நிரல் (niral)}]{\ta{நிரல்} (niral) — a computer program}

\label{lesson:TA-C518-niral}



\begin{quote}
Before the new one: say the Tamil for to disconnect, then the Tamil for to glow.
\end{quote}

\subsection*{You'll want to know: \ta{நிரல்}}
\textbf{\ta{நிரல்}} — \emph{niral} — \textquotedblleft{}a computer program\textquotedblright{}.

\begin{quote}
\ta{என்} \ta{தம்பி} \ta{கணினியில்} \ta{ஒரு} \ta{சிறிய} \ta{நிரல்} \ta{எழுதினான்}.
\end{quote}

\begin{quote}
\emph{eṉ tampi kaṇiṉiyil oru ciṟiya niral eḻutiṉāṉ.}
\end{quote}

\begin{quote}
My younger brother wrote a small program on the computer.
\end{quote}

The one who writes them is a \ta{நிரலாளர்}, a programmer.

\begin{cousinweb}
\textbf{\ta{நிரல்}} (\emph{niral}) is native Dravidian: DEDR 3673 gives \ta{நிரல்} (\emph{niral}), \textquotedblleft{}to be placed in a row, in order\textquotedblright{}; a program is a row of steps in order.
\end{cousinweb}

\begin{grammarlens}[title={What you've built}]
One more word for describing what you lived and felt.
\end{grammarlens}

\subsection*{Your turn}
\begin{itemize}
\raggedright
  \item \emph{Say it:} \emph{niral}
  \item \emph{Say it:} \emph{niral}, once more
  \item \emph{Say it:} \emph{eṉ tampi kaṇiṉiyil oru ciṟiya niral eḻutiṉāṉ.} — My younger brother wrote a small program on the computer.
  \item \emph{Recall:} say the Tamil for to disconnect, then the Tamil for to glow, then say \emph{niral} again
\end{itemize}

\begin{tcolorbox}[breakable,colback=teal!4,colframe=teal!35!black,title={Before you move on}]
What does \ta{நிரல்} mean? (\textquotedblleft{}A computer program\textquotedblright{}.) Say it once more.
\end{tcolorbox}
\section[\texorpdfstring{\ta{குறியீடு} (kuṟiyīṭu)}{குறியீடு (kuṟiyīṭu)}]{\ta{குறியீடு} (kuṟiyīṭu) — a code; a symbol}

\label{lesson:TA-C518-kuriyitu}



\begin{quote}
Before the new one: say the Tamil for to glow, then the Tamil for a computer program.
\end{quote}

\subsection*{You'll want to know: \ta{குறியீடு}}
\textbf{\ta{குறியீடு}} — \emph{kuṟiyīṭu} — \textquotedblleft{}a code; a symbol\textquotedblright{}.

\begin{quote}
\ta{இந்தக்} \ta{கதவைத்} \ta{திறக்க} \ta{ஒரு} \ta{ரகசியக்} \ta{குறியீடு} \ta{வேண்டும்}.
\end{quote}

\begin{quote}
\emph{intak katavait tiṟakka oru rakaciyak kuṟiyīṭu vēṇṭum.}
\end{quote}

\begin{quote}
To open this door you need a secret code.
\end{quote}

\ta{குறி} \textquotedblleft{}a mark\textquotedblright{} with the noun ending \emph{-īṭu}: a mark that stands for something.

\begin{cousinweb}
\textbf{\ta{குறியீடு}} (\emph{kuṟiyīṭu}) is native Dravidian: DEDR 1847 gives \ta{குறி} (\emph{kuṟi}), \textquotedblleft{}a mark, a sign, a symbol\textquotedblright{}; the ending \emph{-īṭu} makes a noun of the act of marking.
\end{cousinweb}

\begin{grammarlens}[title={What you've built}]
One more word for describing what you lived and felt.
\end{grammarlens}

\subsection*{Your turn}
\begin{itemize}
\raggedright
  \item \emph{Say it:} \emph{kuṟiyīṭu}
  \item \emph{Say it:} \emph{kuṟiyīṭu}, once more
  \item \emph{Say it:} \emph{intak katavait tiṟakka oru rakaciyak kuṟiyīṭu vēṇṭum.} — To open this door you need a secret code.
  \item \emph{Recall:} say the Tamil for to glow, then the Tamil for a computer program, then say \emph{kuṟiyīṭu} again
\end{itemize}

\begin{tcolorbox}[breakable,colback=teal!4,colframe=teal!35!black,title={Before you move on}]
What does \ta{குறியீடு} mean? (\textquotedblleft{}A code; a symbol\textquotedblright{}.) Say it once more.
\end{tcolorbox}
\section[\texorpdfstring{\ta{காப்புப்பிரதி} (kāppuppirati)}{காப்புப்பிரதி (kāppuppirati)}]{\ta{காப்புப்பிரதி} (kāppuppirati) — a backup copy}

\label{lesson:TA-C518-kappuppirati}



\begin{quote}
Before the new one: say the Tamil for a computer program, then the Tamil for a code.
\end{quote}

\subsection*{You'll want to know: \ta{காப்புப்பிரதி}}
\textbf{\ta{காப்புப்பிரதி}} — \emph{kāppuppirati} — \textquotedblleft{}a backup copy\textquotedblright{}.

\begin{quote}
\ta{அவள்} \ta{எல்லாப்} \ta{படங்களுக்கும்} \ta{காப்புப்பிரதி} \ta{எடுத்தாள்}.
\end{quote}

\begin{quote}
\emph{avaḷ ellāp paṭaṅkaḷukkum kāppuppirati eṭuttāḷ.}
\end{quote}

\begin{quote}
She made a backup of all the photos.
\end{quote}

One makes one with \ta{எடு}: \ta{காப்புப்பிரதி} \ta{எடு}.

\begin{cousinweb}
\textbf{\ta{காப்புப்பிரதி}} (\emph{kāppuppirati}) is a modern compound: the native \ta{காப்பு} (\emph{kāppu}), \textquotedblleft{}protection\textquotedblright{}, from \ta{கா} (\emph{kā}), \textquotedblleft{}to guard\textquotedblright{} (DEDR 1416), and \ta{பிரதி} (\emph{pirati}), \textquotedblleft{}a copy\textquotedblright{}, from Sanskrit \emph{prati}, \textquotedblleft{}against, in return\textquotedblright{} (Monier-Williams).
\end{cousinweb}

\begin{grammarlens}[title={What you've built}]
One more word for describing what you lived and felt.
\end{grammarlens}

\subsection*{Your turn}
\begin{itemize}
\raggedright
  \item \emph{Say it:} \emph{kāppuppirati}
  \item \emph{Say it:} \emph{kāppuppirati}, once more
  \item \emph{Say it:} \emph{avaḷ ellāp paṭaṅkaḷukkum kāppuppirati eṭuttāḷ.} — She made a backup of all the photos.
  \item \emph{Recall:} say the Tamil for a computer program, then the Tamil for a code, then say \emph{kāppuppirati} again
\end{itemize}

\begin{tcolorbox}[breakable,colback=teal!4,colframe=teal!35!black,title={Before you move on}]
What does \ta{காப்புப்பிரதி} mean? (\textquotedblleft{}A backup copy\textquotedblright{}.) Say it once more.
\end{tcolorbox}
\section[\texorpdfstring{\ta{ரோபோ} (rōpō)}{ரோபோ (rōpō)}]{\ta{ரோபோ} (rōpō) — a robot}

\label{lesson:TA-C518-ropo}



\begin{quote}
Before the new one: say the Tamil for a code, then the Tamil for a backup copy.
\end{quote}

\subsection*{You'll want to know: \ta{ரோபோ}}
\textbf{\ta{ரோபோ}} — \emph{rōpō} — \textquotedblleft{}a robot\textquotedblright{}.

\begin{quote}
\ta{அந்தத்} \ta{தொழிற்சாலையில்} \ta{ரோபோக்கள்} \ta{கார்களைச்} \ta{செய்கின்றன}.
\end{quote}

\begin{quote}
\emph{antat toḻiṟcālaiyil rōpōkkaḷ kārkaḷaic ceykiṉṟaṉa.}
\end{quote}

\begin{quote}
In that factory robots make cars.
\end{quote}

Plural: \ta{ரோபோக்கள்}, \emph{rōpōkkaḷ}.

Its plural doubles the stop: \ta{ரோபோக்கள்}.

\begin{cousinweb}
\textbf{\ta{ரோபோ}} (\emph{rōpō}) is a loan from English \emph{robot}; Tamil drops the final \emph{t}.
\end{cousinweb}

\begin{grammarlens}[title={What you've built}]
One more word for describing what you lived and felt.
\end{grammarlens}

\subsection*{Your turn}
\begin{itemize}
\raggedright
  \item \emph{Say it:} \emph{rōpō}
  \item \emph{Say it:} \emph{rōpō}, once more
  \item \emph{Say it:} \emph{antat toḻiṟcālaiyil rōpōkkaḷ kārkaḷaic ceykiṉṟaṉa.} — In that factory robots make cars.
  \item \emph{Recall:} say the Tamil for a code, then the Tamil for a backup copy, then say \emph{rōpō} again
\end{itemize}

\begin{tcolorbox}[breakable,colback=teal!4,colframe=teal!35!black,title={Before you move on}]
What does \ta{ரோபோ} mean? (\textquotedblleft{}A robot\textquotedblright{}.) Say it once more.
\end{tcolorbox}
\section[\texorpdfstring{\ta{செயற்கை} (ceyaṟkai)}{செயற்கை (ceyaṟkai)}]{\ta{செயற்கை} (ceyaṟkai) — artificial; man-made}

\label{lesson:TA-C518-ceyarkai}



\begin{quote}
Before the new one: say the Tamil for a backup copy, then the Tamil for a robot.
\end{quote}

\subsection*{You'll want to know: \ta{செயற்கை}}
\textbf{\ta{செயற்கை}} — \emph{ceyaṟkai} — \textquotedblleft{}artificial; man-made\textquotedblright{}.

\begin{quote}
\ta{இந்தப்} \ta{பூக்கள்} \ta{செயற்கை}, \ta{உண்மையானவை} \ta{அல்ல}.
\end{quote}

\begin{quote}
\emph{intap pūkkaḷ ceyaṟkai, uṇmaiyāṉavai alla.}
\end{quote}

\begin{quote}
These flowers are artificial, not real.
\end{quote}

It stands before a noun: \ta{செயற்கை} \ta{ஏரி}, an artificial lake. Its opposite is \ta{இயற்கை} (\emph{iyaṟkai}), nature.

\begin{cousinweb}
\textbf{\ta{செயற்கை}} (\emph{ceyaṟkai}) is native Dravidian: DEDR 1957 lists \ta{செயற்கை} (\emph{ceyaṟkai}), \textquotedblleft{}action\textquotedblright{}, under \ta{செய்} (\emph{cey}), \textquotedblleft{}to do, to make\textquotedblright{}. Today it means \textquotedblleft{}man-made\textquotedblright{}, the pair of \ta{இயற்கை}, nature.
\end{cousinweb}

\begin{grammarlens}[title={What you've built}]
One more word for describing what you lived and felt.
\end{grammarlens}

\subsection*{Your turn}
\begin{itemize}
\raggedright
  \item \emph{Say it:} \emph{ceyaṟkai}
  \item \emph{Say it:} \emph{ceyaṟkai}, once more
  \item \emph{Say it:} \emph{intap pūkkaḷ ceyaṟkai, uṇmaiyāṉavai alla.} — These flowers are artificial, not real.
  \item \emph{Recall:} say the Tamil for a backup copy, then the Tamil for a robot, then say \emph{ceyaṟkai} again
\end{itemize}

\begin{tcolorbox}[breakable,colback=teal!4,colframe=teal!35!black,title={Before you move on}]
What does \ta{செயற்கை} mean? (\textquotedblleft{}Artificial; man-made\textquotedblright{}.) Say it once more.
\end{tcolorbox}
